% =====================================================================
% CHAPTER 3: Engineering Design and Methodology
% New department template (2026-07-14), restructure batch B1.
% Merges four overlapping sources into ONE coherent chapter:
%   old mainmatter/ch3.tex  (design, blocks, equations, indices)
%   old ch1 1.4 Research Methodology
%   old ch2 2.6 Detailed Review of the Adopted Methodology
%   old ch7 7.2 (methodology framing only) + 7.5 Review and Feedback
% RED headings: 3.1 Introduction (first), Summary (last).
% Every \label preserved; equations byte-identical; numbers frozen.
% Hard chapter numbers replaced by role-based references per contract.
% =====================================================================
\chapter{Engineering Design and Methodology}
\label{ch:design}

\section{Introduction}
This chapter presents the engineering design of StereoLite, the compact
stereo matching network developed in this thesis, together with the
methodology that produced it. We open with the evidence-gated research
methodology that governed every design decision, then state the hardware
envelope, the design criteria, and the controlled steps by which each
candidate was accepted or rejected. From there the chapter builds the
system model, traces the conceptual design to its literature ancestry,
records the preliminary design decisions that shaped the final network,
and works block by block through the detailed design with all governing
equations. The chapter closes with the performance indices used to
assess the design and the review process behind it. The training
objective and the edge-inference optimization belong to the
implementation chapter, and the experimental results to the results
chapter.

\FloatBarrier
\section{Research Methodology}
\label{sec:research_methodology}
The work follows an evidence-gated design methodology. We admit an
architectural choice only after it wins a controlled comparison, and we
judge every trained model on accuracy, latency, and memory together
rather than on any single metric. Design proceeds by single-knob
ablation on a fixed harness of one hundred training pairs, so that
competing variants are compared under identical data, seed, and compute.
The surviving configuration is then trained on the full dataset on cloud
GPUs and evaluated in three ways: in domain, zero-shot on an unseen real
benchmark, and under controlled rectification perturbation. It is finally
optimized and measured on the target edge device. The steps are:
\begin{enumerate}
\item Study of the stereo matching literature and its efficient and
edge-deployable families;
\item Definition of the design requirements and the edge deployment
envelope;
\item Architecture design by controlled single-knob ablation on a fixed
100-pair harness;
\item Full-dataset training of the selected configuration on cloud GPUs;
\item Evaluation in domain, zero-shot on an unseen real domain, and under
rectification perturbation;
\item Optimization of the network for edge inference and on-device
latency measurement.
\end{enumerate}

The architecture itself was not invented from a blank page. From the
surveyed accuracy frontier we adopted a small set of mechanisms, each
chosen because the surveyed evidence shows it delivers a disproportionate
share of its parent method's accuracy at a cost the edge envelope can
absorb. The conceptual design later in this chapter details how each
mechanism was adapted; the point here is that adoption was itself gated
on evidence, not on preference. This division of labor, evidence for
selection and controlled comparison for tuning, is what keeps the design
auditable decision by decision.

Figure \ref{fig:methodflow} summarizes this pipeline end to end,
including the two feedback paths: a rejected ablation arm returns to the
harness, and a failed evaluation reopens the design itself.

\begin{figure}[!htbp]
\centering
\includegraphics[width=\textwidth]{fig_3_12_methodflow_preview}
\caption{Evidence-gated design methodology from literature through ablation, training, and evaluation to edge deployment.}
\label{fig:methodflow}
\end{figure}

\FloatBarrier
\section{Design Requirements and Requisites}
\label{sec:design_requirements}
The design requirements follow from the class of edge accelerators found
on mobile robots and drones. Together they form a hard envelope, and we
reject a candidate that violates any one of them regardless of accuracy.
The network must stay under three million parameters, so that it is
compact enough for on-device memory and storage; parameter count is also
the first-order proxy for a deployed model's download and update
footprint. It must run in real time on the target edge accelerator, an
NVIDIA Jetson Orin Nano, which delivers 20 to 40 tera-operations per
second of eight-bit throughput \cite{nvidia_orin_nano}, so a forward pass
must fit inside a live camera stream's frame budget on that device. Its
inference working set at the deployment precision must stay under 200
megabytes, small enough to coexist with the rest of a robot's software
stack on a shared-memory edge board, though a higher-precision
development GPU may use more. Finally, every operator on the inference
path must quantize cleanly to eight-bit integer arithmetic, since the
accelerators in this class reach their rated throughput only under
integer execution.

\FloatBarrier
\section{Design Criteria}
Within the envelope above, we judge candidate designs on balance rather
than on any single number. Every design decision is evaluated on the full
eight-metric stereo set: end-point error, root-mean-square error, median
absolute error, the bad-0.5, bad-1, bad-2, and bad-3 rates, and the
D1-all rate, together with inference latency and peak memory. These
quantities are defined formally in Section \ref{sec:design_assessment}. A
variant that improves one metric while degrading the others, or that buys
accuracy with latency the envelope cannot afford, is not accepted.
Because in-domain gains can mask a loss of generalization, we require a
cross-dataset zero-shot evaluation after every training run. All
comparisons are made under equal conditions, the same training pairs,
random seed, and step budget, so that a measured difference is
attributable to the design change alone.

\FloatBarrier
\section{Design Steps}
The architectural choices in this chapter were made under a controlled
matched-overfit protocol rather than intuition. We train each candidate
on a fixed set of SceneFlow pairs with a fixed seed, deterministic
convolution kernels, and an identical step budget and evaluation
protocol, and we report all eight stereo metrics plus an inference
latency benchmark for every arm. Single knobs are varied one at a time,
the winning criterion is registered before a run so the outcome cannot be
rationalized after the fact, and every candidate is checked against a
held-out cross-dataset evaluation, so that an in-domain gain which fails
to transfer is not mistaken for progress. This section states the
protocol; the runs it produced appear in the results chapter, and the
mapping from each decision to its supporting run is collected in the
preliminary design below.

\FloatBarrier
\section{Modeling of Proposed System}
Figure \ref{fig:architecture} gives the complete pipeline. StereoLite is
a coarse-to-fine tile-hypothesis network with light recurrent refinement.
A shared feature encoder processes both views and a separate left-image
context encoder feeds the recurrent state; a group-wise cost volume at
the coarsest scale initializes a slanted tile state; convolutional GRU
updates refine the state across three scales with plane-aware propagation
between them; a narrow-band geometry encoding volume corrects the
estimate through a learned fail-soft gate; and plane rendering with
learned convex upsampling recovers the full-resolution disparity.

% Rotated full-page float: the architecture overview is ~4:1 wide, so
% at \textwidth it printed ~5 pt text. Sideways gives it the full page
% height (~1.5x larger) without redesigning the diagram.
\begin{sidewaysfigure}
\centering
% INTERIM: viewer render of figures/fig_3_1_architecture.drawio. Once
% hand-adjustments in draw.io are done, export a cropped PDF over
% fig_3_1_architecture.pdf and swap the filename below.
\includegraphics[width=\textheight]{fig_3_1_architecture_preview}
\caption{StereoLite architecture: five numbered stages with supervised outputs marked by red dots.}
\label{fig:architecture}
\end{sidewaysfigure}

Table \ref{tab:scales} summarizes the spatial sizes and the main
operation at each pipeline scale for the default 384 by 640 input.

\begin{table}[!htbp]
\centering
\caption{Spatial size, channel count, and main operation at each pipeline scale.}
\label{tab:scales}
\small
\begin{tabular}{lccl}
\toprule
Scale & Spatial size & Channels & Main operation\\
\midrule
Full & $384\times 640$ & 3 & input and output\\
1/2 & $192\times 320$ & 32 & convex-upsample guidance\\
1/4 & $96\times 160$ & 128 & GEV fusion and 3 refinements\\
1/8 & $48\times 80$ & 256 & 3 refinements\\
1/16 & $24\times 40$ & 256 & cost initialization and 2 refinements\\
\bottomrule
\end{tabular}
\end{table}

\FloatBarrier
\section{Modeling of Conceptual Design}
\label{sec:conceptual_design}
Conceptually, the pipeline is an assembly of five stages, each of which
adapts a mechanism from the literature review to the edge budget. We did
not adopt these mechanisms because they are recent, but because the
surveyed evidence shows each one carries a large share of its parent
method's accuracy at a cost our envelope can afford. The detailed design
of each stage follows in Section \ref{sec:detailed_design}; here we set
out the reasoning behind each choice.

The matching cost is the group-wise correlation volume of GwcNet
\cite{guo2019gwcnet}, given in \eqref{eq:gwc_lit}. It carries more
matching evidence than a plain correlation yet costs a small fraction of
a concatenation volume, and it degrades gracefully when the
three-dimensional aggregation behind it is cut to a few layers, which is
exactly the regime a compact network operates in. We restrict it to a
single coarse scale, which keeps the peak activation memory within the
small inference working set the envelope demands.

The refinement operator and the final upsampling come as a pair from
RAFT-Stereo \cite{lipson2021raftstereo}, given in \eqref{eq:convgru_lit}:
a convolutional gated recurrent unit reused every iteration, and a
learned convex combination that restores full resolution from a coarse
estimate with image-derived weights instead of a dense fine-scale volume.
The recurrent update concentrates capacity in a small module, so accuracy
scales with compute rather than with parameters. We run the pair for only
a handful of iterations at coarse resolution, which fits the real-time
budget that rules out the many-iteration desktop operating points.

The third mechanism is the geometry encoding volume of IGEV-Stereo
\cite{xu2023igev}, given in \eqref{eq:gev_lit}. Its warm-start disparity
saves the early iterations that a cold-started recurrent loop wastes,
which is the single largest lever for cutting iteration count, and the
surveyed ablations show a handful of warm-started iterations matching
dozens of cold ones. In our design a learned gate blends this global
proposal into the recurrent estimate, so the volume corrects rather than
competes.

The fourth mechanism is the slanted tile-plane representation of HITNet
\cite{tankovich2021hitnet}. A grid of tiles, each carrying a disparity,
two plane slopes, and a descriptor, replaces the dense cost volume as the
working state of the network, and the slopes preserve sub-pixel structure
across scales without rebuilding a volume at every level. This is the
most aggressive of the surveyed cost-volume compressions, and its
published operating point, under one million parameters at real-time
latency, sits squarely inside our envelope.

The fifth choice is structural rather than borrowed from a single method:
the matching path and the recurrent context are kept in separate streams.
A shared feature encoder feeds the matching cost, and a dedicated
left-image context encoder feeds the recurrent state. If the matching
features also served as context, the recurrent unit could not separate
appearance from correspondence; the RAFT-Stereo lineage
\cite{lipson2021raftstereo} established that a dedicated context path
improves refinement, and we adopt the same separation. What the surveyed
literature does not directly address is the specific combination these
mechanisms form, a tile-hypothesis network of a few million parameters
that fuses a geometry encoding volume through a fail-soft gate, measured
for near real-time inference on the Jetson Orin Nano class of hardware
and characterized for robustness to imperfect rectification. That gap is
the design problem this thesis takes up.

\FloatBarrier
\section{Preliminary Design}
The final architecture was reached through a sequence of preliminary
variants, each accepted or rejected under the design-steps protocol
above. Every load-bearing decision in this chapter is tied to a
controlled study in the results chapter or to published evidence in the
literature review. Table \ref{tab:design_evidence} collects this mapping
so the design can be audited decision by decision; it is the record of
preliminary designs, including the rejected ones, that shaped the final
system.

\begin{table}[H]
\centering
\caption{Architectural and training decisions with their supporting evidence.}
\label{tab:design_evidence}
\small
\begin{tabular}{p{5.2cm}p{9.0cm}}
\toprule
Decision & Evidence\\
\midrule
YOLO26s shared encoder & encoder ablation: best EPE of three encoders at +2 ms\\
Separate left context encoder & RAFT-Stereo lineage \cite{lipson2021raftstereo}; context-gated variants win the family ablation\\
Group-wise cost at one coarse scale & GwcNet \cite{guo2019gwcnet} graceful degradation under small aggregation budgets\\
Iterated coarse-to-fine refinement & refinement ablation: single-pass collapses, finer scales add cost without gain\\
Local $\pm 2$ correlation lookup & CREStereo \cite{li2022crestereo}: local search beats all-pairs; also confers rectification tolerance\\
Slanted tile planes & HITNet \cite{tankovich2021hitnet}; plane rendering wins the blur study\\
Narrow GEV band + fail-soft gate & family ablation best EPE; narrowing accuracy-neutral\\
Rejected: selective recurrent unit & family ablation: EPE worsens on the compact chassis\\
Rejected: fresh 1/4 initialization & family ablation: EPE and bad-1 worsen\\
Rejected: fix compositions & blur study: worse than plane rendering alone\\
Threshold-stack + D1 objective & loss ablation: wins EPE, bad-2, D1 together\\
Native-crop input protocol & input-protocol study: native axis EPE 6.67 to 2.87 at no resized-axis cost\\
Augmentation triplet & matched comparison: best val EPE 2.778 to 1.921\\
OneCycle 60k schedule & training dynamics: validation floor reached before the budget ends\\
\bottomrule
\end{tabular}
\end{table}

\FloatBarrier
\section{Detailed Design}
\label{sec:detailed_design}
This section works through each pipeline block in the order data flows
through it, together with the governing equations for each stage. Read
at a high level before the detail, the pipeline does five things in
sequence. It first turns the two raw images into features a network can
compare (Section~\ref{sec:detailed_design}.1). It then makes a coarse,
low-resolution guess at the disparity of every image tile by testing a
range of candidate horizontal shifts and keeping the one whose left and
right features agree best (\ref{sec:detailed_design}.2). That guess is
usually wrong in the details, so a recurrent loop repeatedly looks at how
well the current guess lines up the two views and nudges it closer,
scale by scale, carrying the estimate from coarse tiles to progressively
finer ones as it goes (\ref{sec:detailed_design}.3-\ref{sec:detailed_design}.4).
Because that loop only ever looks at a small neighborhood around each
tile, a second, wider-reaching check is added purely to correct the
tiles the loop is still struggling with, without disturbing the ones it
has already resolved well (\ref{sec:detailed_design}.5). Finally the
coarse tile grid is rendered and sharpened up to the full image
resolution (\ref{sec:detailed_design}.6), and the resulting disparity map
is converted into an actual three-dimensional point cloud
(\ref{sec:detailed_design}.7). Figure~\ref{fig:architecture} already
showed this same flow as a numbered block diagram; the subsections below
give each block's mechanism and governing equations in turn.

Figure~\ref{fig:progression} makes the same process concrete by running
the trained network on a real stereo pair and showing the disparity map
as it actually leaves each stage. The coarse initialization at one
sixteenth resolution is blocky and only roughly right; each later stage
sharpens it, until the full-resolution output resolves the motorcycle's
wheels, frame, and background boxes cleanly. The panel labelled
\emph{GEV proposal} and the fail-soft gate beside the inputs preview the
correction step of Section~\ref{sec:detailed_design}.5, where a second,
independent estimate is blended in only where the gate opens. Reading the
subsections below alongside this figure shows both what each block does
and what it contributes to the final disparity.

\begin{figure}[!htbp]
\centering
\includegraphics[width=\textwidth]{fig_3_9_progression}
\caption[Disparity as it leaves each stage of the pipeline, on a real
Middlebury 2014 pair.]{Disparity as it leaves each stage of the trained
network, run on a real Middlebury 2014 pair (Motorcycle
\cite{scharstein2014middlebury}). Every panel is an actual model output,
upsampled to full resolution and shown on one shared disparity scale;
the estimate sharpens from the coarse one-sixteenth initialization to the
full-resolution result. The gate $w$ (top right) shows where the GEV
proposal is trusted.}
\label{fig:progression}
\end{figure}

\subsection{Shared Feature Encoder and Context Encoder}
Before the two views can be compared, each has to be turned into
features a network can match, and the left view alone has to supply the
running context the recurrent refinement in
Section~\ref{sec:detailed_design}.3 will need. This subsection builds
those two encoders: one shared between the left and right image so both
are described in the same feature space, and one private to the left
image alone.

The matching network is a truncated YOLO26s backbone whose first seven
blocks are retained, producing features at one quarter, one eighth, and
one sixteenth resolution with 128, 256, and 256 channels. The choice of a
detection backbone is deliberate: its features are pretrained to localize
objects across scales, which is close to what stereo matching needs from
an encoder, and the encoder ablation confirms that this backbone beats
both a hand-built GhostConv stack and its own smaller sibling on every
accuracy metric at a two-millisecond cost. We remove the detection neck and head above the one-sixteenth stage
at construction rather than leaving them unused, because dead blocks
would waste parameters and forward compute. The two views are
concatenated along the batch dimension and pass through a single shared
invocation, so both images are processed with identical weights. A
separate context encoder processes the left image alone: a strided
convolution is followed by two GhostConv blocks with
squeeze-and-excitation, and the result is projected to a 32-channel
quarter-resolution context map that is pooled to the coarser scales. We
keep matching evidence and recurrent context in separate streams for a
specific reason: if the matching features also served as context, the
recurrent unit could not separate appearance from correspondence. The
RAFT-Stereo lineage \cite{lipson2021raftstereo} established that a
dedicated context path improves refinement, and we adopt the same
separation.

\begin{figure}[H]
\centering
\includegraphics[width=\textwidth]{fig_3_2_encoders_preview}
\caption{Shared YOLO26s matching encoder and separate left-image context encoder.}
\label{fig:encoders}
\end{figure}

\subsection{Tile Initialization and Cost Volume}
This subsection produces the network's first, coarse guess at disparity.
The idea is the same one every stereo method relies on: for a small
patch of the left image, try a range of candidate horizontal shifts,
compare each shifted patch of the right image against it, and keep the
shift whose features agree best. Doing this once, cheaply, at a coarse
resolution gives every tile a starting point that later stages only need
to correct rather than discover from nothing.

The tile initialization module builds a group-wise correlation volume,
following GwcNet \cite{guo2019gwcnet}, at one sixteenth resolution. For
each of eight feature groups and each
integer disparity, the matching score is the mean channel-wise product
between the left features and the horizontally shifted right features, as
given in \eqref{eq:gwc}. A small three-dimensional convolutional network
with an eight, sixteen, sixteen, one channel schedule regularizes the
volume, and a softmax along the disparity axis yields a probability
distribution. The initial disparity and confidence come from the
GC-Net-style soft-argmin \cite{kendall2017gcnet} and the distribution
peak, as in \eqref{eq:softargmax_init}.
The plane slopes are initialized to zero and the hidden feature is
replaced by the coarsest context map. The 24 hypotheses at this scale
span roughly 0 to 368 full-resolution pixels.

\begin{equation}
C(g,d,y,x) \;=\; \frac{1}{C/G}\sum_{c=1}^{C/G}
  f_{L}^{\,g,c}(y,x)\; f_{R}^{\,g,c}(y,\,x-d)
\label{eq:gwc}
\end{equation}
where $G = 8$ is the number of feature groups, $C$ the channel count,
$d \in \{0, \dots, 23\}$ the integer disparity hypotheses at one sixteenth
resolution, and $f_{L}^{\,g,c}$, $f_{R}^{\,g,c}$ the $c$-th channel of
group $g$ of the left and right feature maps.

\begin{equation}
d_{0} \;=\; \sum_{d} p(d)\, d, \qquad
c_{0} \;=\; \max_{d}\, p(d)
\label{eq:softargmax_init}
\end{equation}
where $p(d)$ is the softmax of the regularized cost volume along the
disparity dimension, $d_0$ the expected disparity, and $c_0$ the peak
probability that serves as an initial confidence.

\begin{figure}[!htbp]
\centering
\includegraphics[width=0.95\textwidth]{fig_3_3_tile_init_preview}
\caption{Tile-state initialization: group-wise correlation volume regularized by 3-D CNN with soft-argmax outputs. The 24 hypotheses at 1/16 scale cover roughly 0 to 368 full-resolution pixels.}
\label{fig:tile_init}
\end{figure}

\subsection{Recurrent Tile Refinement}
The initial guess from Section~\ref{sec:detailed_design}.2 is coarse and
often wrong in the details, so this subsection corrects it iteratively
rather than in one shot. At each step, the network checks how well the
current disparity estimate lines up the two views by warping one onto
the other and looking at the local mismatch, then nudges the estimate,
the local surface tilt, and its own confidence a little closer to
correct. Repeating this update, and carrying a persistent memory of
what has already been tried between repetitions, lets the network spend
more computation on tiles that are hard to match and less on tiles it
already got right, rather than paying the same fixed cost everywhere.

Separate refinement modules refine the tile state at one sixteenth, one
eighth, and one quarter resolution; their weights are not shared, because
the matching-feature channel counts differ across scales. At each
iteration the right feature map is warped toward the left view by the
current disparity, and a differentiable local correlation, in the manner
of CREStereo's local search \cite{li2022crestereo}, is evaluated
over five offsets within a half-range of two pixels. The recurrent input
assembles the left features, the warped right features, the current
state, the local correlation, and the context, following the
context-conditioned input RAFT-Stereo feeds to its recurrent unit
\cite{lipson2021raftstereo}, as in
\eqref{eq:gru_input}. A convolutional gated recurrent unit, following
RAFT-Stereo \cite{lipson2021raftstereo}, then updates
the persistent hidden state through the gate equations of
\eqref{eq:convgru}. A two-layer head with 48 hidden channels predicts
residuals for disparity, both plane slopes, and confidence, and a second
learned gate emits four independent sigmoid weights that scale each
residual. A softplus constrains disparity to remain non-negative. The
schedule runs two iterations at one sixteenth, three at one eighth, and
three at one quarter, for eight learned updates in total.

\begin{equation}
x \;=\; \bigl[\, f_{L},\; \mathcal{W}(f_{R};\,d),\; d,\; s_{x},\; s_{y},\;
  c,\; \mathrm{corr},\; \mathrm{ctx} \,\bigr]
\label{eq:gru_input}
\end{equation}
where $\mathcal{W}(f_{R}; d)$ denotes the right features warped by the
current disparity, $\mathrm{corr}$ the five-offset local correlation, and
$\mathrm{ctx}$ the scale-matched context map.

\begin{equation}
\begin{aligned}
z &= \sigma\!\bigl(\mathrm{Conv}([h,\,x])\bigr), &
r &= \sigma\!\bigl(\mathrm{Conv}([h,\,x])\bigr),\\
q &= \tanh\!\bigl(\mathrm{Conv}([r \odot h,\,x])\bigr), &
h' &= (1-z) \odot h + z \odot q,
\end{aligned}
\label{eq:convgru}
\end{equation}
where $h$ is the 32-channel persistent hidden state, $z$ and $r$ the
update and reset gates, $q$ the candidate state, and $\odot$ the
element-wise product.

\begin{figure}[!htbp]
\centering
\includegraphics[width=\textwidth]{fig_3_4_refinement_preview}
\caption{One recurrent refinement iteration with warping, local correlation lookup, and gated residual updates.}
\label{fig:refinement}
\end{figure}

\subsection{Plane-Aware Cross-Scale Propagation}
Refinement at one resolution has to hand its result up to the next,
finer resolution before refinement can continue there, and a generic
upscaling would just blur the coarse tile boundaries across the new,
smaller tiles. This subsection is the hand-over step, and it does
better than plain upscaling because it also carries forward the local
surface tilt each tile already estimated, so a slanted surface is placed
correctly at the finer grid instead of being flattened and then
re-discovered from scratch.

Between adjacent scales, the propagation module resamples the disparity,
slopes, hidden feature, and confidence to the finer grid. Disparity is
multiplied by two because its pixel unit changes with resolution. Unlike
plain bilinear interpolation, the propagated disparity also uses the
local plane slopes to place the four child pixels around each parent
tile, as given in \eqref{eq:plane_prop}. This tile-plane component, in
the manner of HITNet \cite{tankovich2021hitnet}, preserves a local
slanted-surface hypothesis while moving to a finer grid; sub-pixel
structure is retained without rebuilding a cost volume at every level.

\begin{equation}
d_{\mathrm{child}} \;=\; 2\,\mathcal{B}(d_{\mathrm{parent}})
  \;+\; 2\,s_{x}\,\Delta x \;+\; 2\,s_{y}\,\Delta y,
\qquad \Delta x, \Delta y \in \{-0.25,\, +0.25\}
\label{eq:plane_prop}
\end{equation}
where $\mathcal{B}$ is bilinear resampling of the parent disparity, $s_x$
and $s_y$ the local plane slopes, and $\Delta x, \Delta y \in \{-0.25,
+0.25\}$ the child offsets; the factor of two rescales the disparity to
the finer grid.

\begin{figure}[!htbp]
\centering
\includegraphics[width=\textwidth]{fig_3_5_plane_prop_preview}
\caption{Plane-aware cross-scale propagation using local slopes instead of bilinear interpolation. The slopes shift each child off the parent plane, and the factor of two rescales the disparity to the finer grid.}
\label{fig:plane_prop}
\end{figure}

\subsection{Geometry Encoding Volume and Fail-Soft Fusion}
The recurrent loop in Section~\ref{sec:detailed_design}.3 is a local
process: every update only looks at a small neighborhood around the
current estimate, so a tile whose initial guess started far from the
truth can stay stuck near that bad guess no matter how many iterations
run. This subsection adds a second, independent estimate with a wider
reach, used only to pull those stuck tiles toward a better value. It is
deliberately not allowed to overrule tiles the recurrent loop has
already resolved well: a learned gate decides, tile by tile, how much
weight this second opinion gets, and it learns to give it almost none
where the recurrent estimate is already trustworthy.

The geometry encoding volume adds an independent global proposal at one
quarter resolution, in the spirit of IGEV-Stereo \cite{xu2023igev}. In
the trained configuration it is narrow: rather than the full 64 disparity
planes of the base design, it samples 33 hypotheses in a band of plus or
minus 16 around the current tile disparity. The efficiency study
validates this change as accuracy-neutral, and it also removes roughly a
quarter of the model's multiply-accumulate operations. The band is less restrictive than it sounds: 16
quarter-resolution planes correspond to 64 full-resolution pixels on
either side of the current estimate, and since the coarse initialization
has already localized the disparity, the band's job is correction, not
search. The full-range volume spends most of its computation scoring
hypotheses the tile state has already ruled out, so removing them costs
nothing. The volume is regularized by three three-dimensional
convolutions with 16 hidden channels, and the softmax distribution yields
the expected disparity, a confidence, and a projected geometry feature,
following IGEV-Stereo's warm-start soft-argmin \cite{xu2023igev}, as in
\eqref{eq:gev_expect}. The proposal is not substituted
unconditionally: a learned gate takes the context, the geometry feature,
both confidences, the absolute disagreement, and the current disparity,
and produces a scalar weight through \eqref{eq:failsoft}. We initialize
the gate's final bias to minus four, so the initial weight is
approximately 0.018 and training begins close to the stable recurrent
baseline; the network then learns where the geometry proposal helps. This
fail-soft design stops the global branch from corrupting regions the
recurrent path already resolves well.

\begin{equation}
d_{\mathrm{gev}} = \sum_{d} p(d)\, d, \qquad
c_{\mathrm{gev}} = \max_{d}\, p(d), \qquad
g_{\mathrm{gev}} = \phi\Bigl(\sum_{d} G(d)\, p(d)\Bigr)
\label{eq:gev_expect}
\end{equation}
where $p(d)$ is the softmax over the regularized quarter-resolution
volume, $G(d)$ the regularized three-dimensional geometry feature, and
$\phi$ a projection to a 16-channel two-dimensional expected-geometry map.

\begin{equation}
\begin{aligned}
w &= \sigma\!\Bigl(F\bigl([\mathrm{ctx}_{4},\, g_{\mathrm{gev}},\, c,\,
  c_{\mathrm{gev}},\, \lvert d_{\mathrm{gev}} - d \rvert,\, d]\bigr)\Bigr),\\
d_{\mathrm{fused}} &= \mathrm{softplus}\bigl(d + w\,(d_{\mathrm{gev}} - d)\bigr),
\end{aligned}
\label{eq:failsoft}
\end{equation}
where $F$ is the fusion network whose final-layer bias is initialized to
$-4$, $c$ and $c_{\mathrm{gev}}$ the tile and geometry confidences, and
$w$ the blend weight; the slopes are attenuated by $(1 - w)$ and the
confidence becomes the maximum of the two estimates.

\begin{figure}[!htbp]
\centering
\includegraphics[width=\textwidth]{fig_3_6_gev_fusion_preview}
\caption{Narrow-band geometry encoding volume fused into tile estimates via learned fail-soft gate.}
\label{fig:gev_fusion}
\end{figure}

\subsection{Plane Rendering and Learned Convex Upsampling}
Every stage so far has worked on a tile grid coarser than the final
output, at one quarter of the image resolution at best, so the last step
is to turn that coarse grid into a full-resolution disparity map. A
naive upsampling would blur object edges across tile boundaries; this
subsection instead uses the slanted-plane estimate each tile already
carries to reconstruct a sharp surface within it, and a learned weighting
scheme, rather than fixed bilinear interpolation, to combine neighboring
tiles so that boundaries stay crisp where the scene actually has a depth
discontinuity.

The final one-quarter disparity is upsampled to full resolution in two
learned two-times stages. The first stage predicts masks from the
one-quarter left features and the second from the one-half left features.
For every output subpixel, a mask predicts nine weights over the
three-by-three coarse neighborhood, and a softmax makes the weights
non-negative and sum to one, so the upsampled disparity is a convex
combination of nearby coarse values. This is the same convex upsampling
used in RAFT-Stereo \cite{lipson2021raftstereo}. Before upsampling, a
plane-rendering step reconstructs each tile from its disparity and slopes
and applies an edge gate; this was the decisive fix in the
boundary-sharpness study, where it reduced the bad-0.5 rate by roughly 22
percent relative to the control.
Convex upsampling can select different source neighbors near a boundary,
but it cannot invent a disparity hypothesis absent from the one-quarter
prediction, so thin and distant structures remain a limitation.

\begin{figure}[!htbp]
\centering
\includegraphics[width=0.92\textwidth]{fig_3_7_upsample_preview}
\caption{Learned convex upsampling via softmax-weighted 3x3 neighborhoods to full resolution. Two 2$\times$ stages (1/4 to 1/2 to full) each multiply disparity by two, with non-negative weights that sum to one.}
\label{fig:upsample}
\end{figure}

\subsection{From Disparity to 3D Reconstruction}
\label{sec:reconstruction_method}
The disparity map itself is an intermediate result; the goal of the
pipeline is the three-dimensional structure of the scene that the
disparity encodes. Once the disparity is known, every valid pixel
back-projects to a metric three-dimensional point through the pinhole
relations
\begin{equation}
Z = \frac{f\,B}{d}, \qquad
X = \frac{(u - c_x)\,Z}{f}, \qquad
Y = \frac{(v - c_y)\,Z}{f}
\label{eq:backprojection}
\end{equation}
where $(u, v)$ are the pixel coordinates, $(c_x, c_y)$ the principal
point, $f$ the focal length in pixels, $B$ the stereo baseline, and
$(X, Y, Z)$ the recovered point in the camera frame. Each point is
colored with the corresponding left-image pixel, which yields a dense
colored point cloud of the scene. For the reconstructions reported in the
results chapter, the cloud is post-processed with two standard
operations: a statistical outlier filter removes the isolated flying
points that appear along depth discontinuities, and a voxel downsampling
makes the point density uniform. The stereo rig's calibrated values, a
focal length of approximately 1005 pixels at the native 1280-pixel width
and a baseline of 52 millimeters, scale to the inference resolution by the
width ratio. This reconstruction path is what the thesis title refers to:
the network provides the dense disparity, and the calibrated geometry
turns it into a metric three-dimensional model of the scene.

\FloatBarrier
\section{Design Assessment and Performance Indices}
\label{sec:design_assessment}
We report accuracy with the standard stereo metrics in
\eqref{eq:metrics}: end-point error is the mean absolute disparity error
over valid pixels, bad-$t$ is the percentage of valid pixels whose
absolute error exceeds $t$ pixels, and D1-all is the KITTI outlier
definition \cite{menze2015kitti}. We also report root-mean-square error and median absolute
error: the first is sensitive to large outliers, while the second gives
an outlier-robust summary. Each image contributes its own averaged
metric, so scores are averaged per image rather than pooled across every
pixel. Latency is measured at batch one, 384 by 640 resolution, with ten
warm-up and a hundred timed passes; memory is the peak inference working
set, and the device used is stated throughout.

\begin{equation}
\begin{aligned}
\mathrm{EPE} &= \frac{1}{\lvert V \rvert} \sum_{p \in V}
  \bigl\lvert \hat{d}(p) - d^{*}(p) \bigr\rvert, \qquad
\mathrm{bad}\text{-}t = \frac{100}{\lvert V \rvert}\,
  \Bigl\lvert \bigl\{ p \in V :
  \lvert \hat{d}(p) - d^{*}(p) \rvert > t \bigr\} \Bigr\rvert,\\
\mathrm{D1} &= \frac{100}{\lvert V \rvert}\, \Bigl\lvert \bigl\{ p \in V :
  \lvert \hat{d}(p) - d^{*}(p) \rvert > 3 \;\wedge\;
  \lvert \hat{d}(p) - d^{*}(p) \rvert > 0.05\, d^{*}(p) \bigr\} \Bigr\rvert
\end{aligned}
\label{eq:metrics}
\end{equation}
where $V$ is the set of valid pixels, $\hat{d}$ the predicted disparity,
$d^{*}$ the ground truth, and $t \in \{0.5, 1, 2, 3\}$ pixels the bad-rate
thresholds.

Alongside accuracy, latency, and memory, the parameter budget is itself a
performance index, bounded by the envelope in Section
\ref{sec:design_requirements}. Figure \ref{fig:param_budget} reports the
module-wise split of the trained model; the split is computed directly
from the checkpoint's state dictionary.

\begin{figure}[!htbp]
\centering
\includegraphics[width=0.8\textwidth]{fig_3_11_param_budget}
\caption{Module-wise parameter distribution: encoder and refinement stages hold 95\% of 2.96 M parameters.}
\label{fig:param_budget}
\end{figure}

These indices did not merely score the finished model; they governed the
design as it was built. We followed a review-and-feedback process rather
than unchecked iteration: every proposed change, whether architectural, a
loss term, or an input protocol, had to win a controlled ablation on the
matched harness described above before entering the production
configuration, and a change that lost its comparison was not adopted no
matter how plausible it looked. Every run, win or lose, went into a master
experiments log with its full configuration and metrics, so the evidence
behind each decision is auditable and no comparison depends on memory. We
reviewed training curves, qualitative panels, and ablation tables with the
supervisor at regular intervals, and that feedback redirected the work
more than once, most visibly toward the cross-dataset evaluation gate that
now follows every training run. Failed variants were documented rather
than discarded: the rejected mechanisms and anti-synergistic compositions
recorded in the results chapter constrain future work as firmly as the
adopted ones justify the present design.

\FloatBarrier
\section{Summary}
This chapter set out the engineering design of StereoLite and the
methodology that produced it. We work under an evidence-gated
methodology: an architectural choice is admitted only after it wins a
controlled single-knob comparison on a fixed harness, and every model is
judged on accuracy, latency, and memory together with a mandatory
cross-dataset zero-shot evaluation. The design obeys a hard envelope:
under 3 million parameters, real-time operation on a 4 to 6 TOPS
accelerator, under 200 megabytes of inference memory, and
INT8-exportable operators. Within that envelope we assembled StereoLite
from five mechanisms adapted from the surveyed literature, a group-wise
cost volume, recurrent tile refinement with plane-aware cross-scale
propagation, a narrow-band geometry encoding volume with a fail-soft
gate, and plane rendering with learned convex upsampling, fed by a shared
matching encoder and a separate context encoder. Each block was specified
with its governing equations, and each design decision is auditable
through the evidence table and the master experiments log.

% SOURCES:
%   old mainmatter/ch3.tex: 3.1 Introduction, 3.3 Design Requirements
%     (sec:design_requirements), 3.4 Design Criteria, 3.5 Design Steps,
%     3.6 Modeling of Proposed System (fig:architecture, tab:scales),
%     3.7 Modeling of Conceptual Design, 3.8 Preliminary Design
%     (tab:design_evidence), 3.9 Detailed Design (sec:detailed_design;
%     all blocks + eq:gwc, eq:softargmax_init, eq:gru_input, eq:convgru,
%     eq:plane_prop, eq:gev_expect, eq:failsoft, eq:backprojection;
%     figs 3.2-3.7; sec:reconstruction_method), 3.10 Design Assessment
%     (sec:design_assessment; eq:metrics, fig:param_budget), 3.11 Summary.
%   old ch1.tex 1.4 Research Methodology -> 3.2 Research Methodology
%     (evidence-gated 6-step enumerate).
%   old ch2.tex 2.6 Detailed Review of the Adopted Methodology -> the
%     four-mechanism selection rationale folded into 3.2 (bridge) and 3.7
%     Modeling of Conceptual Design; \eqref{eq:gwc_lit},
%     \eqref{eq:convgru_lit}, \eqref{eq:gev_lit} referenced (defined in ch2).
%   old ch7.tex 7.5 Review and Feedback -> folded into 3.10 Design
%     Assessment (master log, supervisor review, documented failures).
% DROPPED:
%   old ch3 3.2 Problem Statements + design assumptions: problem
%     statements/RQs moved to ch1 (see MOVED-AWAY); design assumptions
%     (rectified input, 192 px max disparity, synthetic GT) not re-stated
%     to avoid duplicating ch1's problem framing.
%   old ch2 2.6 opening trade-off framing with tab:method_comparison and
%     fig:pareto: those are comparative-analysis artifacts owned by ch2
%     2.3; not duplicated here (see MOVED-AWAY). Kept only the
%     four-mechanism adoption prose, which is design rationale.
%   old ch3 3.7 compact 5-stage prose: superseded by the richer
%     per-mechanism edge-justification merged from ch2 2.6 (better version
%     won); no unique content lost.
%   Hard chapter-number literals ("Chapter 2/7/8") replaced with
%     role-based references ("literature review", "implementation
%     chapter", "results chapter") per contract rule 5; specific ablation
%     tables now cited via \ref to their ch6 labels.
% MOVED-AWAY:
%   old ch3 3.2 Problem Statements + four research questions -> ch1 1.3
%     Research Questions / Problem Statements (owner ch1).
%   old ch7 7.2 subsection Training Objective (eq:loss, fig:supervision)
%     -> implementation chapter (ch4); NOT reproduced here to avoid label
%     duplication. Only 7.2's methodology framing overlaps design steps.
%   old ch7 7.2 subsection Datasets and Training Protocol
%     (fig:input_protocol, tab:train_config) -> data collection (ch5).
%   old ch2 2.6 tab:method_comparison and fig:pareto -> ch2 2.3
%     Comparative Analysis (owner ch2); referenced by role only.
